\documentclass[10pt,a4paper]{amsart}
\usepackage{amsmath,amssymb,amsthm,graphicx}
\usepackage{verbatim}
\newenvironment{reference}{\comment}{\endcomment}
%\newenvironment{reference}{}{}
\newenvironment{slogan}{\comment}{\endcomment}
\newenvironment{history}{\comment}{\endcomment}
% For commutative diagrams we use Xy-pic
\usepackage[all]{xy}

% We use 2cell for 2-commutative diagrams.
\xyoption{2cell}
% Theorem environments.
%
\theoremstyle{plain}
\newtheorem{theorem}[subsection]{Theorem}
\newtheorem{proposition}[subsection]{Proposition}
\newtheorem{lemma}[subsection]{Lemma}

\theoremstyle{definition}
\newtheorem{definition}[subsection]{Definition}
\newtheorem{example}[subsection]{Example}
\newtheorem{exercise}[subsection]{Exercise}
\newtheorem{situation}[subsection]{Situation}

\theoremstyle{remark}
\newtheorem{remark}[subsection]{Remark}
\newtheorem{remarks}[subsection]{Remarks}

\numberwithin{equation}{subsection}

% Macros
%
\def\lim{\mathop{\mathrm{lim}}\nolimits}
\def\colim{\mathop{\mathrm{colim}}\nolimits}
\def\Spec{\mathop{\mathrm{Spec}}}
\def\Hom{\mathop{\mathrm{Hom}}\nolimits}
\def\Ext{\mathop{\mathrm{Ext}}\nolimits}
\def\SheafHom{\mathop{\mathcal{H}\!\mathit{om}}\nolimits}
\def\SheafExt{\mathop{\mathcal{E}\!\mathit{xt}}\nolimits}
\def\Sch{\mathit{Sch}}
\def\Mor{\mathop{\mathrm{Mor}}\nolimits}
\def\Ob{\mathop{\mathrm{Ob}}\nolimits}
\def\Sh{\mathop{\mathit{Sh}}\nolimits}
\def\NL{\mathop{N\!L}\nolimits}
\def\CH{\mathop{\mathrm{CH}}\nolimits}
\def\proetale{{pro\text{-}\acute{e}tale}}
\def\etale{{\acute{e}tale}}
\def\QCoh{\mathit{QCoh}}
\def\Ker{\mathop{\mathrm{Ker}}}
\def\Im{\mathop{\mathrm{Im}}}
\def\Coker{\mathop{\mathrm{Coker}}}
\def\Coim{\mathop{\mathrm{Coim}}}

% Boxtimes
%
\DeclareMathSymbol{\boxtimes}{\mathbin}{AMSa}{"02}

%
% Macros for moduli stacks/spaces
%
\def\QCohstack{\mathcal{QC}\!\mathit{oh}}
\def\Cohstack{\mathcal{C}\!\mathit{oh}}
\def\Spacesstack{\mathcal{S}\!\mathit{paces}}
\def\Quotfunctor{\mathrm{Quot}}
\def\Hilbfunctor{\mathrm{Hilb}}
\def\Curvesstack{\mathcal{C}\!\mathit{urves}}
\def\Polarizedstack{\mathcal{P}\!\mathit{olarized}}
\def\Complexesstack{\mathcal{C}\!\mathit{omplexes}}
% \Pic is the operator that assigns to X its picard group, usage \Pic(X)
% \Picardstack_{X/B} denotes the Picard stack of X over B
% \Picardfunctor_{X/B} denotes the Picard functor of X over B
\def\Pic{\mathop{\mathrm{Pic}}\nolimits}
\def\Picardstack{\mathcal{P}\!\mathit{ic}}
\def\Picardfunctor{\mathrm{Pic}}
\def\Deformationcategory{\mathcal{D}\!\mathit{ef}}
\usepackage{cleveref}
\crefname{originallemma}{original lemma}{original lemmas}
\Crefname{originallemma}{Original lemma}{Original lemmas}
\makeatletter
\newlabel{lemma-log-complex}{{0FMW}{source}{}{sourcepointer-1}{}}
\newlabel{lemma-log-complex@cref}{{[originallemma][0][]0FMW}{[1][0][]source}{}{}{}}
\makeatother
\crefname{originallemma}{original lemma}{original lemmas}
\Crefname{originallemma}{Original lemma}{Original lemmas}
\makeatletter
\newlabel{lemma-the-complex-for-L-star}{{0FUF}{source}{}{sourcepointer-2}{}}
\newlabel{lemma-the-complex-for-L-star@cref}{{[originallemma][0][]0FUF}{[1][0][]source}{}{}{}}
\makeatother
\crefname{originallemma}{original lemma}{original lemmas}
\Crefname{originallemma}{Original lemma}{Original lemmas}
\makeatletter
\newlabel{lemma-the-complex-for-L-star-gives-chern-class}{{0FUG}{source}{}{sourcepointer-3}{}}
\newlabel{lemma-the-complex-for-L-star-gives-chern-class@cref}{{[originallemma][0][]0FUG}{[1][0][]source}{}{}{}}
\makeatother
\crefname{originalremark}{original remark}{original remarks}
\Crefname{originalremark}{Original remark}{Original remarks}
\makeatletter
\newlabel{remark-cup-product-as-a-map}{{0FU7}{source}{}{sourcepointer-4}{}}
\newlabel{remark-cup-product-as-a-map@cref}{{[originalremark][0][]0FU7}{[1][0][]source}{}{}{}}
\makeatother
\title{The Gysin type boundary of the log complex is compatible with the maps attached to the first Chern class of the divisor}
\author{The Stacks Project Authors}
\date{stacks-project:derham.tex:local-snapshot-20261012; GNU FDL 1.2 or later}
\begin{document}
\maketitle
\noindent\textit{Editorial note.} All delivered mathematical text is the author's own, copied from the pinned Stacks Project chapter file derham.tex; no statement or proof step was rewritten, reordered or removed.\par\medskip
\noindent\textit{Editorial note.} Preamble provenance: the pinned chapter file opens with \textbackslash{}input\{preamble\} and carries no preamble text of its own, so the project's single Stacks-wide preamble staged at sources/stacks-derham/preamble.tex is excerpted verbatim as three spans, lines 8-12 (verbatim package and the project's own reference/slogan/history display environments), lines 14-18 (Xy-pic setup) and lines 154-217 (theorem declarations and mathematical macros). The wrapper class and the amsmath/amssymb/amsthm/graphicx base packages are supplied by the assembler. No macro was invented or redefined.\par\medskip
\noindent\textit{Editorial note.} Same-chapter pointers are located at their \textbackslash{}label lines inside the pinned chapter file; cross-chapter pointers are located in the project's own tag list staged in the snapshot (sources/stacks-chow/tags.txt), at the verbatim lines carrying their canonical tags.\par\medskip
\noindent\textit{Allowed standard prerequisites (editorial).}\begin{itemize}
\item The short exact sequence defining the de Rham complex of log poles and its boundary. Same chapter derham.tex, label lemma-log-complex (invoked at proof lines 2568, 2635 and 2672). It gives the shifted sequence \$0 \textbackslash{}to \textbackslash{}Omega\textasciicircum{}\textbackslash{}bullet\_\{X/S\}[1] \textbackslash{}to \textbackslash{}Omega\textasciicircum{}\textbackslash{}bullet\_\{X/S\}(\textbackslash{}log Y)[1] \textbackslash{}to \textbackslash{}Omega\textasciicircum{}\textbackslash{}bullet\_\{Y/S\} \textbackslash{}to 0\$ whose boundary is the map \$\textbackslash{}delta\$ of the statement.
\item The map of the derived category attached to a de Rham class. Same chapter derham.tex, remark remark-cup-product-as-a-map (invoked at proof lines 2573 and 2579). It defines \$\textbackslash{}xi'\$ and \$\textbackslash{}zeta'\$ as the maps attached to \$c\_1\textasciicircum{}\{dR\}(\textbackslash{}mathcal\{O\}\_X(-Y))\$ and to its restriction to \$Y\$.
\item The punctured line bundle sequence and the identification of its boundary with the Chern class map. Same chapter derham.tex, labels lemma-the-complex-for-L-star (invoked at proof lines 2608 and 2643) and lemma-the-complex-for-L-star-gives-chern-class (invoked at proof lines 2606 and 2641). They describe \$\textbackslash{}xi'\$ and \$\textbackslash{}zeta'\$ as boundary maps of the punctured line bundle sequences, which reduces the commutativity to producing morphisms of short exact sequences.
\end{itemize}\medskip
\ifdefined\hypertarget\hypertarget{sourcepointer-1}{}\fi
\noindent\textit{Source reference.} Original lemma 0FMW: sources/stacks-derham/pinned-original-source.tex, source lines 2391--2391. Navigation pointer for \textbackslash{}ref\{lemma-log-complex\} used inside the counted span (line 2568); the argument that consumes it is retained in the span.\par\medskip
\ifdefined\hypertarget\hypertarget{sourcepointer-2}{}\fi
\noindent\textit{Source reference.} Original lemma 0FUF: sources/stacks-derham/pinned-original-source.tex, source lines 1348--1348. Navigation pointer for \textbackslash{}ref\{lemma-the-complex-for-L-star\} used inside the counted span (line 2608); the argument that consumes it is retained in the span.\par\medskip
\ifdefined\hypertarget\hypertarget{sourcepointer-3}{}\fi
\noindent\textit{Source reference.} Original lemma 0FUG: sources/stacks-derham/pinned-original-source.tex, source lines 1422--1422. Navigation pointer for \textbackslash{}ref\{lemma-the-complex-for-L-star-gives-chern-class\} used inside the counted span (line 2606); the argument that consumes it is retained in the span.\par\medskip
\ifdefined\hypertarget\hypertarget{sourcepointer-4}{}\fi
\noindent\textit{Source reference.} Original remark 0FU7: sources/stacks-derham/pinned-original-source.tex, source lines 575--575. Navigation pointer for \textbackslash{}ref\{remark-cup-product-as-a-map\} used inside the counted span (line 2573); the argument that consumes it is retained in the span.\par\medskip
\begin{lemma}
\label{lemma-gysin-via-log-complex}
Let $f : X \to S$ be a morphism of schemes. Let $Y \subset X$ be an effective
Cartier divisor. Assume the de Rham complex of log poles is defined for
$Y \subset X$ over $S$. Denote
$$
\delta : \Omega^\bullet_{Y/S} \to \Omega^\bullet_{X/S}[2]
$$
in $D(X, f^{-1}\mathcal{O}_S)$ the ``boundary'' map coming from the
short exact sequence in Lemma \ref{lemma-log-complex}. Denote
$$
\xi' : \Omega^\bullet_{X/S} \to \Omega^\bullet_{X/S}[2]
$$
in $D(X, f^{-1}\mathcal{O}_S)$ the map of
Remark \ref{remark-cup-product-as-a-map}
corresponding to $\xi = c_1^{dR}(\mathcal{O}_X(-Y))$. Denote
$$
\zeta' : \Omega^\bullet_{Y/S} \to \Omega^\bullet_{Y/S}[2]
$$
in $D(Y, f|_Y^{-1}\mathcal{O}_S)$ the map of
Remark \ref{remark-cup-product-as-a-map} corresponding to
$\zeta = c_1^{dR}(\mathcal{O}_X(-Y)|_Y)$. Then the diagram
$$
\xymatrix{
\Omega^\bullet_{X/S} \ar[d]_{\xi'} \ar[r] &
\Omega^\bullet_{Y/S} \ar[d]^{\zeta'} \ar[ld]_\delta \\
\Omega^\bullet_{X/S}[2] \ar[r] &
\Omega^\bullet_{Y/S}[2]
}
$$
is commutative in $D(X, f^{-1}\mathcal{O}_S)$.
\end{lemma}

\begin{proof}
More precisely, we define $\delta$ as the boundary map corresponding to the
shifted short exact sequence
$$
0 \to \Omega^\bullet_{X/S}[1] \to
\Omega^\bullet_{X/S}(\log Y)[1] \to
\Omega^\bullet_{Y/S} \to 0
$$
It suffices to prove each triangle commutes. Set
$\mathcal{L} = \mathcal{O}_X(-Y)$. Denote $\pi : L \to X$ the line bundle
with $\pi_*\mathcal{O}_L = \bigoplus_{n \geq 0} \mathcal{L}^{\otimes n}$.

\medskip\noindent
Commutativity of the upper left triangle.
By Lemma \ref{lemma-the-complex-for-L-star-gives-chern-class}
the map $\xi'$ is the boundary map of the triangle given in
Lemma \ref{lemma-the-complex-for-L-star}.
By functoriality it suffices to prove there exists a morphism of
short exact sequences
$$
\xymatrix{
0 \ar[r] &
\Omega^\bullet_{X/S}[1] \ar[r] \ar[d] &
\Omega^\bullet_{L^\star/S, 0}[1] \ar[r] \ar[d] &
\Omega^\bullet_{X/S} \ar[r] \ar[d] &
0 \\
0 \ar[r] &
\Omega^\bullet_{X/S}[1] \ar[r] &
\Omega^\bullet_{X/S}(\log Y)[1] \ar[r] &
\Omega^\bullet_{Y/S} \ar[r] &
0
}
$$
where the left and right vertical arrows are the obvious ones.
We can define the middle vertical arrow by the rule
$$
\omega' + \text{d}\log(s) \wedge \omega \longmapsto
\omega' + \text{d}\log(f) \wedge \omega
$$
where $\omega', \omega$ are local sections of $\Omega^\bullet_{X/S}$
and where $s$ is a local generator of $\mathcal{L}$ and
$f \in \mathcal{O}_X(-Y)$ is the corresponding section of the ideal
sheaf of $Y$ in $X$. Since the constructions of the maps in
Lemmas \ref{lemma-the-complex-for-L-star} and \ref{lemma-log-complex}
match exactly, this works.

\medskip\noindent
Commutativity of the lower right triangle. Denote
$\overline{L}$ the restriction of $L$ to $Y$.
By Lemma \ref{lemma-the-complex-for-L-star-gives-chern-class}
the map $\zeta'$ is the boundary map of the triangle given in
Lemma \ref{lemma-the-complex-for-L-star} using the line bundle
$\overline{L}$ on $Y$.
By functoriality it suffices to prove there exists a morphism of
short exact sequences
$$
\xymatrix{
0 \ar[r] &
\Omega^\bullet_{X/S}[1] \ar[r] \ar[d] &
\Omega^\bullet_{X/S}(\log Y)[1] \ar[r] \ar[d] &
\Omega^\bullet_{Y/S} \ar[r] \ar[d] &
0 \\
0 \ar[r] &
\Omega^\bullet_{Y/S}[1] \ar[r] &
\Omega^\bullet_{\overline{L}^\star/S, 0}[1] \ar[r] &
\Omega^\bullet_{Y/S} \ar[r] &
0 \\
}
$$
where the left and right vertical arrows are the obvious ones.
We can define the middle vertical arrow by the rule
$$
\omega' + \text{d}\log(f) \wedge \omega \longmapsto
\omega'|_Y + \text{d}\log(\overline{s}) \wedge \omega|_Y
$$
where $\omega', \omega$ are local sections of $\Omega^\bullet_{X/S}$
and where $f$ is a local generator of $\mathcal{O}_X(-Y)$ viewed as
a function on $X$ and where $\overline{s}$ is $f|_Y$ viewed as a
section of $\mathcal{L}|_Y = \mathcal{O}_X(-Y)|_Y$.
Since the constructions of the maps in
Lemmas \ref{lemma-the-complex-for-L-star} and \ref{lemma-log-complex}
match exactly, this works.
\end{proof}

\end{document}
